\DocumentMetadata{
  pdfversion=2.0,
  pdfstandard=ua-2,
  lang=en-US,
  tagging=on,
  tagging-setup={math/setup=mathml-SE,tabsorder=structure}
}
\documentclass[11pt,letterpaper]{article}
\usepackage[margin=0.68in,includefoot]{geometry}
\usepackage{newcomputermodern}
\usepackage{amsmath,amscd,mathtools}
\usepackage{tikz,adjustbox}
\providecommand{\square}{\mdlgwhtsquare}
\usepackage[hidelinks]{hyperref}
\hypersetup{
  pdftitle={Algebra PhD Exam, September 1988},
  pdfauthor={University of Florida Department of Mathematics},
  pdfsubject={Graduate mathematics examination},
  pdfdisplaydoctitle=true,
  bookmarksopen=true
}
\setcounter{secnumdepth}{0}
\setlength{\parindent}{0pt}
\setlength{\parskip}{0.28em}
\setlength{\emergencystretch}{3em}
\setlength{\leftmargini}{2.15em}
\setlength{\leftmarginii}{2.65em}
\setlength{\leftmarginiii}{3em}
\pagestyle{empty}
\raggedbottom
\allowdisplaybreaks
\NewTaggingSocketPlug{block/list/label}{examproblem}{%
  \tagstructbegin{tag=Lbl}\tagstructend%
  \tagstructbegin{tag=\LBody}%
  \tagstructbegin{tag=\ProblemHeadingTag,title-o={\ProblemHeadingText}}%
  \tagmcbegin{tag=\ProblemHeadingTag,actualtext={\ProblemHeadingText}}%
  #2%
  \tagmcend%
  \tagstructend%
}
\newenvironment{examproblems}[2]{%
  \def\ProblemHeadingTag{#2}%
  \AssignTaggingSocketPlug{block/list/label}{examproblem}%
  \begin{enumerate}%
  \setcounter{enumi}{#1}%
  \renewcommand{\labelenumi}{\textbf{\theenumi.}}%
  \setlength{\topsep}{0.35em}%
  \setlength{\partopsep}{0pt}%
  \setlength{\itemsep}{0.48em}%
  \setlength{\parsep}{0.18em}%
}{\end{enumerate}}
\newenvironment{examsubparts}{%
  \AssignTaggingSocketPlug{block/list/label}{default}%
  \begin{enumerate}%
  \setlength{\topsep}{0.18em}%
  \setlength{\partopsep}{0pt}%
  \setlength{\itemsep}{0.16em}%
  \setlength{\parsep}{0.1em}%
}{\end{enumerate}}
\newenvironment{examinstructions}{%
  \par\begingroup\itshape
}{%
  \par\endgroup\vspace{0.18em}
}
\makeatletter
\renewcommand\subsection{\@startsection{subsection}{2}{0pt}%
  {-1.25ex plus -.25ex minus -.1ex}%
  {0.55ex plus .1ex}%
  {\normalfont\large\bfseries}}
\renewcommand{\@maketitle}{%
  \newpage\null\vskip -1em%
  {\footnotesize\bfseries
    \href{https://www.ufl.edu/}{University of Florida}, \href{https://math.ufl.edu/}{Department of Mathematics}\par}%
  \vskip -0.5em%
  \tagstructbegin{tag=H1,title={Algebra PhD Exam, September 1988}}%
  \tagpdfparaOff%
  \tagmcbegin{tag=H1}%
    {\LARGE\bfseries\href{https://gma.math.ufl.edu/exams/algebra-phd/}{Algebra PhD Exam}, September 1988\par}%
  \tagmcend%
  \tagpdfparaOn%
  \tagstructend%
  \vskip 0.25em%
  \tagmcbegin{artifact}%
    \hrule height 1pt%
  \tagmcend%
  \vskip 1em%
}
\makeatother
\title{Algebra PhD Exam, September 1988}
\author{}
\date{}
\begin{document}
\maketitle
\thispagestyle{empty}
\begin{examinstructions}
Do seven problems, at least two from each section.
\end{examinstructions}
\subsection{I. Group Theory.}
\begin{examproblems}{0}{H3}
\def\ProblemHeadingText{Problem 1.}
\item
Let~\(G\) be the group of all \(3\times 3\) non-singular matrices of determinant~\(1\) with coefficients in the field of~\(3\) elements.
\begin{examsubparts}
\item[\textnormal{(a)}]
Find \(|G|\).
\item[\textnormal{(b)}]
Find a Sylow~\(3\)-subgroup of~\(G\).
\item[\textnormal{(c)}]
Compute \(|\operatorname{Syl}_3(G)|\), where \(\operatorname{Syl}_3(G)\) is the set of all Sylow~\(3\)-subgroups of~\(G\).
\end{examsubparts}
\def\ProblemHeadingText{Problem 2.}
\item
Let~\(G\) be a finite group acting transitively on a set~\(X\), \(x\in X\), and~\(P\) be a Sylow~\(p\)-subgroup of \(G_x\) for some prime~\(p\). Consider the set \(\operatorname{Fix}(P)=\{y\in X\mid y^P=y\}\), the set of fixed points of~\(P\).
\begin{examsubparts}
\item[\textnormal{(a)}]
Prove that \(N_G(P)\) acts on \(\operatorname{Fix}(P)\), i.e., maps \(\operatorname{Fix}(P)\) to \(\operatorname{Fix}(P)\).
\item[\textnormal{(b)}]
Prove that \(N_G(P)\) acts transitively on \(\operatorname{Fix}(P)\).
\end{examsubparts}
\def\ProblemHeadingText{Problem 3.}
\item
Two permutation representations \(G\xrightarrow{\pi_i}\Sigma(X_i)\), (the symmetric group on \(X_i\) are equivalent if there exists \(\alpha:X_1\to X_2\) such that for all \(g\in G\), \(x\in X_1\), we have \(x(g\pi_1)\alpha=(x\alpha)(g\pi_2)\). Prove that the following two permutation representations of \(S_4\) of degree \(12\) are not equivalent:
\begin{examsubparts}
\item[\textnormal{(a)}]
on the right coset space of \(\langle(12)\rangle\).
\item[\textnormal{(b)}]
on the right coset space of \(\langle(12)(34)\rangle\).
\end{examsubparts}
\def\ProblemHeadingText{Problem 4.}
\item
Let~\(G\) be a finite group and \(A\leq \operatorname{Aut}(G)\) such that \(G=[G,A]\). Suppose \(N\leq G\) satisfying \([N,A]=1\). Use the 3-subgroup lemma to prove \(N\leq Z(G)\), the center of~\(G\). (Recall the 3-subgroups lemma: For three subgroups \(X,Y,Z\) of a group~\(M\), if \([X,Y,Z]=1\) and \([Y,Z,X]=1\), then \([Z,X,Y]=1\).)
\end{examproblems}
\subsection{II. Homological algebra.}
\begin{examproblems}{0}{H3}
\def\ProblemHeadingText{Problem 1.}
\item
Let~\(R\) be a ring and consider the following commutative diagram of~\(R\)-modules and~\(R\)-module homomorphisms such that each row is a short exact sequence.
{\tagpdfsetup{math/mathml/structelem=false,math/mathml/AF=false,math/alt/use=true}
\[
\begin{CD}
0 @>>> A @>{f}>> B @>{g}>> C @>>> 0 \\
@. @V{\alpha}VV @V{\beta}VV @V{\gamma}VV @. \\
0 @>>> A' @>{f'}>> B' @>{g'}>> C' @>>> 0
\end{CD}
\]
}
 Prove that if~\(\alpha\) and~\(\gamma\) are isomorphisms then~\(\beta\) is also an isomorphism.
\def\ProblemHeadingText{Problem 2.}
\item
Prove the \(\mathbb Z\)-module isomorphisms, where \(c=(m,n)\) is the greatest common divisor.
\begin{examsubparts}
\item[\textnormal{(a)}]
\(\mathbb Z/m\mathbb Z\otimes\mathbb Z/n\mathbb Z=\mathbb Z/c\mathbb Z\).
\item[\textnormal{(b)}]
\(\operatorname{Hom}(\mathbb Z/m\mathbb Z,\mathbb Z/n\mathbb Z)=\mathbb Z/c\mathbb Z\).
\end{examsubparts}
\def\ProblemHeadingText{Problem 3.}
\item
A module~\(P\) over a ring~\(R\) is said to be projective if, given any diagram of~\(R\)-module homomorphisms with~\(g\) surjective
\[
\begin{array}{ccc}
& & P \\
& & \downarrow^{f} \\
A & \xrightarrow{g} & B
\end{array}
\]
 there exists an~\(R\)-module homomorphism \(h:P\longrightarrow A\) such that the following diagram is commutative
\[
\begin{array}{ccc}
& & P \\
& \overset{h}{\swarrow} & \downarrow^{f} \\
A & \xrightarrow{g} & B
\end{array}
\]
 Prove that the following conditions on a ring~\(R\) with identity are equivalent:
\begin{examsubparts}
\item[\textnormal{(a)}]
Every~\(R\)-module is projective.
\item[\textnormal{(b)}]
Every short exact sequence of~\(R\)-modules splits.
\end{examsubparts}
\end{examproblems}
\subsection{III. Commutative algebra.}
\begin{examproblems}{0}{H3}
\def\ProblemHeadingText{Problem 1.}
\item
A non-empty subset~\(S\) of a ring~\(R\) with identity is called multiplicative if \(ab\in S\) whenever \(a,b\in S\).
\begin{examsubparts}
\item[\textnormal{(a)}]
If~\(I\) is a proper ideal of~\(R\) disjoint from~\(S\), then prove that there exists an ideal~\(P\) maximal with respect to containing~\(I\) and being disjoint from~\(S\). Furthermore, show that~\(P\) is a prime ideal.
\item[\textnormal{(b)}]
If~\(P\) is a prime ideal of ring~\(R\), define the localization of~\(R\) at~\(P\) and prove that the localization has a unique maximal ideal.
\end{examsubparts}
\def\ProblemHeadingText{Problem 2.}
\item
Find a reduced primary decomposition for the following ideals. Also give the prime ideals to which the primary ideals belong.
\begin{examsubparts}
\item[\textnormal{(a)}]
\((24)\) as an ideal in \(\mathbb Z\).
\item[\textnormal{(b)}]
\((2x,x^2)\) as an ideal in \(\mathbb Z[x]\).
\item[\textnormal{(c)}]
What theorems allow us to conclude that every ideal in \(\mathbb Z[x]\) has a primary decomposition.
\end{examsubparts}
\def\ProblemHeadingText{Problem 3.}
\item
Let~\(R\) be a unique factorization domain and~\(K\) its field of fractions.
\begin{examsubparts}
\item[\textnormal{(a)}]
Show that if \(a\in K\) is integral over~\(R\), then \(a\in R\). (In other words, \(R\) is integrally closed.)
\item[\textnormal{(b)}]
True or false: If~\(S\) is a subring of~\(K\) containing~\(R\), then~\(S\) is also a unique factorization domain. If true, prove your answer. If false give examples of~\(R\), \(K\), and~\(S\) that show it not to be true.
\end{examsubparts}
\def\ProblemHeadingText{Problem 4.}
\item
Determine, up to isomorphism, all semisimple rings of order \(1008=2^5\cdot 3^2\cdot 7\).
\end{examproblems}
\end{document}
